% 第 14 章 Creating Routes（续）

\section{在布线上添加缓冲器}
\subsection*{Adding Buffers on Routes}

Add Buffer on Route 工具可以沿指定线网的布线添加缓冲器（buffer）。

要启动 Add Buffer on Route 工具，单击 Command search 按钮，并在搜索框中输入 \cmd{add\_buffer\_on\_route}。

下图中，1 是命令搜索图标，2 是输入命令的文本框，3 显示与搜索文本匹配的命令。

Add Buffer on Route 对话框随即打开，分为以下部分：
\begin{itemize}
  \item 顶部区域（Top region）
  \item Basic 选项卡
  \item Advanced 选项卡
\end{itemize}

对话框中以红色显示、或前面带星号的字段为必填。尚未提供的必填信息会显示在该对话框底部。

表~\ref{tab:gui-buf-route} 列出 Add Buffer on Route 对话框的选项。

\begin{longtable}{@{}p{0.32\textwidth}p{0.62\textwidth}@{}}
\caption{Add Buffer on Route 对话框选项}
\label{tab:gui-buf-route}\\
\toprule
字段 & 说明 \\
\midrule
\endfirsthead
\caption{Add Buffer on Route 对话框选项（续）}\\
\toprule
字段 & 说明 \\
\midrule
\endhead
\midrule
\multicolumn{2}{r}{\small 下页继续} \\
\endfoot
\bottomrule
\endlastfoot
Nets & 指定要插入缓冲器的线网。 \\
Location & 指定缓冲器或反相器对（inverter pair）单元所在位置的中心点坐标。 \\
Specified buffers & 指定缓冲器或反相器对单元的中心点坐标，以及它们的名称和库单元。 \\
Specified bus buffers &
指定总线缓冲器（bus buffer）或反相器对单元的中心点坐标，以及它们的名称和库单元。格式为 \cmd{\{bus\_buffer\_pattern1 lib\_cell1 x1 y1 bus\_buffer\_pattern2 lib\_cell2 x2 y2 ...\}}。 \\
Repeater distance & 指定沿布线插入的各个缓冲器之间的距离，单位为用户单位（user unit）。 \\
\multicolumn{2}{@{}l}{\textbf{Mode}} \\
\addlinespace
Repeater distance length ratio &
以总布线长度的比例，指定沿布线插入的各个缓冲器之间的距离。 \\
Point & 指定版图视图上的点坐标。 \\
Layer & 从下拉列表中选择一层。 \\
Detect layer & 指定工具自动检测位置坐标所在的层。 \\
Name & 指定所添加缓冲器的名称。 \\
Lib Cell & 指定用作中继器（repeater）的库单元对象。 \\
Distance & 指定中继器距离，即各个缓冲器之间的距离。 \\
Scale & 按该层指定的缩放因子，缩放缓冲器之间的距离。 \\
Ratio & 各个缓冲器之间的距离与总布线长度之比。 \\
Scaled by width &
按该层默认宽度与实际布线宽度之比，缩放缓冲器之间的距离。 \\
\multicolumn{2}{@{}l}{\textbf{Basic 选项卡}} \\
\addlinespace
Only global routed nets & 允许在全局布线的线网（global routed net）上插入缓冲器。 \\
Punch port &
即使拓扑（topology）与逻辑层次中的网表连接不同，也允许在布线段上插入缓冲器。 \\
Respect voltage area &
在电压区域（voltage area）的层次内插入缓冲器，使缓冲器在版图上物理地位于该层次内部。 \\
Reference Cells > Voltage Area, Lib Cell &
指定在给定电压区域中使用的库单元对象。 \\
Respect gas station &
仅在指定供电线网的加油站（gas station）内、并沿布线添加缓冲器。这些加油站的 \cmd{is\_primary} 属性为 false。 \\
Physical feedthrough &
指定可以添加缓冲器的电压区域，即使它们只是物理贯通（physical-only feedthrough）。 \\
Select MV buffers > Voltage Area, Single Rail, Dual Rail &
指定在给定电压区域中使用的单电源轨（single-rail）与双电源轨（dual-rail）库单元对象。这些是多电压缓冲器（multi-voltage buffer，MV buffer）。 \\
Prefix > Bus cell &
指定由 \opt{-user\_specified\_bus\_buffers} 选项添加的新缓冲器所用的前缀。 \\
Prefix > Cell &
指定插入网表的新缓冲器所用的前缀。默认情况下，命令使用 \cmd{eco\_cell} 作为前缀。 \\
Prefix > Net &
指定向网表插入缓冲器时新建线网所用的前缀。默认情况下，命令使用 \cmd{eco\_net} 作为前缀。 \\
\multicolumn{2}{@{}l}{\textbf{Advanced 选项卡}} \\
\addlinespace
Snap to sites & 把新添加的缓冲器吸附到最近的位点（site）。 \\
On top hierarchy &
把缓冲器加到层次最高层的线网上。默认情况下，缓冲器加在其所驱动管脚的最低公共祖先层次中。 \\
Verbose & 显示详细消息。 \\
Respect blockages &
指定不在布局阻挡（placement blockage）、软宏或硬宏内添加缓冲器。 \\
Spare cell aware &
在冻结硅（freeze-silicon）模式下启用感知备用单元（spare cell）的功能。 \\
Insertion over cell > Allowed &
指定可以在其中添加中继器的软宏或硬宏实例。本选项与 Respect blockages 选项一起使用。 \\
Insertion over cell > Not allowed &
指定不允许添加中继器的软宏或硬宏实例。 \\
Max distance > Location to route &
指定给定位置与布线之间的最大水平或垂直距离，单位为用户单位。本选项与 Respect gas station 选项一起使用。启用 Respect gas station 时，缓冲器必须加在加油站内部。 \\
Max distance > Incomplete route & 定义不完整布线的最大搜索距离。 \\
Max distance > Route to gas station &
指定添加缓冲器的加油站到布线的距离。若未指定本选项，加油站必须位于布线上。 \\
Max distance > Spare cell &
在冻结硅模式下，指定新添加的缓冲器与备用单元之间的最大距离。 \\
First distance -- Length &
指定驱动管脚与沿布线引入的第一个缓冲器之间的距离，单位为用户单位。本选项与 Repeater distance 或 Repeater distance length ratio 选项一起使用。 \\
First distance -- Ratio &
以总布线长度的比例，指定驱动管脚与沿布线引入的第一个缓冲器之间的距离。本选项与 Repeater distance length ratio 选项一起使用。 \\
\end{longtable}

表~\ref{tab:gui-gas-station} 给出加油站示例。

\begin{table}[htbp]
\centering
\caption{加油站示例}
\label{tab:gui-gas-station}
\begin{tabularx}{\textwidth}{@{}l X@{}}
\toprule
情形 & 说明 \\
\midrule
启用 Respect gas station &
缓冲器必须加在加油站内部。 \\
Max distance > Route to gas station &
允许放大加油站，因此即使加油站不在布线上，仍可以添加缓冲器。 \\
示例流程 &
\begin{minipage}[t]{\linewidth}
\begin{enumerate}
  \item 选择线网 \cmd{p0}。
  \item 添加用户指定的缓冲器。
\end{enumerate}
\end{minipage} \\
\bottomrule
\end{tabularx}
\end{table}

下图给出结果。

% ============================================================
\section{估算点到点电阻}
\subsection*{Estimating Point-to-Point Resistance}

要计算当前设计中两点（管脚或端口）之间的电阻（resistance），使用 \cmd{estimate\_resistance} 命令。

在该命令中，可以用多种方式指定起点与终点：
\begin{itemize}
  \item 使用 \opt{-net}。本选项只适用于仅有两个管脚或端口、且每个管脚或端口只有一个端子的线网。
  \item 使用 \opt{-pins}。本选项适用于有两个以上管脚或端口的线网。必须在命令中指定确切的管脚和端口，命令才能测量这些指定点之间的电阻。
  \item 使用 \opt{-shapes}。本选项适用于测量复杂结构中任意线段的电阻，例如电源地管脚与网格（mesh）的连接。为兼顾性能，避免把整个网格作为一组形状传入。
\end{itemize}

使用 \opt{-shapes} 时，需指定起点及其所在层，以及终点及其所在层。

用 \opt{-quiet} 与 \opt{-verbose} 分别在控制台上抑制消息，或详细打印消息。

下面的命令估算指定单元上 VDD 管脚之间的电阻：
\begin{lstlisting}
icc2_shell> estimate_resistance -pins "INV_279/VDD U27/VDD"
\end{lstlisting}
